\providecommand{\Perf}{\operatorname{Perf}}
\providecommand{\Fun}{\operatorname{Fun}}
\providecommand{\Locfd}{\operatorname{Loc}_{\mathrm{fd}}}
\providecommand{\alg}{\mathrm{alg}}
\providecommand{\sch}{\mathrm{sch}}
\providecommand{\Qfd}{\mathcal Q_k}
\providecommand{\Sfd}{\mathcal S_k}

\subsection{The missing condition in the proposed correction}

Throughout this part, $k$ is a nonzero field belonging to the universe
of small objects. For a space $X$, a finite-dimensional local system
means a functor from the fundamental groupoid of $X$ to finite-dimensional
$k$-vector spaces, sending arrows to isomorphisms. Ranks may differ on
different components. Write $\Qfd$ for the class of maps $f:X\to Y$ such
that the \emph{specified} pullback map
\[
 H^n(Y;L)\longrightarrow H^n(X;f^*L)
 \quad\text{is an isomorphism for every }L\in\Locfd(Y),\ n\geq0.
 \tag{D.0}
\]
For a functor of small categories we apply this definition to its nerve
realization. Thus coefficients in (D.0) are on the target. There is no
unspecified transport of source coefficients.

\begin{claim}{D.1}{Refuted}{Cohomology with finite-dimensional target coefficients alone is not a localizer}
\label{D1}
For every field $k$, $\Qfd$ fails two-out-of-three, even on functors
between finite, connected, pointed categories. Consequently Part D.1 and
Lead 2, with precisely the proposed cohomology test and no further
condition, are false. Neither characteristic zero nor algebraic closure
nor choosing basepoints repairs this assertion.
\end{claim}
\begin{proof}
Choose a prime $\ell$ different from $\operatorname{char}k$: take
$\ell=2$ unless $\operatorname{char}k=2$, in which case take $\ell=3$.
Let $C_\ell=\mathbb Z/\ell$ and let $\mathbf B C_\ell$ be the category
with one object and endomorphism group $C_\ell$. Let
$i:e\to\mathbf B C_\ell$ send the sole identity to the identity and
$p:\mathbf B C_\ell\to e$ be the unique functor. They satisfy
$pi=\mathrm{id}_e$ in $\Cat$.

The averaging idempotent $\ell^{-1}\sum_{g\in C_\ell}g$ makes the
invariants functor on $k[C_\ell]$-modules exact: average any lift of an
invariant vector to lift it invariantly. Hence all its positive derived
functors vanish. In particular $H^n(BC_\ell;V)$ equals $V$ for $n=0$
and zero for $n>0$ when the action is trivial. These are exactly the
coefficients pulled back from $e$, so $p\in\Qfd$.

For $i$, use the regular representation $L=k[C_\ell]$ on its target.
The map in degree zero is the inclusion
\[
 k[C_\ell]^{C_\ell}\longrightarrow k[C_\ell].
\]
Its source has dimension one and its target has dimension $\ell$.
Thus $i\notin\Qfd$, although $p$ and $pi$ belong to $\Qfd$.
Every category and functor has been specified; the group laws and
coefficient calculation are checked by
\texttt{code/D\_verify\_examples.py}, with output in the corresponding
\texttt{.out} file. Averaging proves the vanishing in all degrees; the
finite chain computation is an independent check through degree four.
\end{proof}

\subsection{Which derived category gives the correct statement}

Let $\Perf(k)$ be the small stable $\infty$-category of perfect
$k$-complexes, and put
\[
 \mathcal P_k(X)=\Fun(X,\Perf(k)),
 \qquad
 \Sfd=\{f:X\to Y:\ f^*:\mathcal P_k(Y)\to\mathcal P_k(X)
                       \text{ is an equivalence}\}.
 \tag{D.1}
\]
Here $X$ denotes the entire fundamental $\infty$-groupoid, rather than
its one-truncation. We use coherent local systems of complexes. On each
connected component, $\mathcal P_k(X)$ is exactly the thick stable
subcategory of $\Fun(X,D(k))$ generated by ordinary finite-dimensional
local systems. On a disconnected space the assertion is componentwise;
no common bound on amplitudes over infinitely many components is imposed.

\begin{claim}{D.2}{Proved}{Full faithfulness and essential surjectivity are different tests}
\label{D2}
For every map of small spaces $f:X\to Y$, the following hold.
\begin{enumerate}
\item $f\in\Qfd$ if and only if $f^*:\mathcal P_k(Y)\to\mathcal P_k(X)$
is fully faithful.
\item $f\in\Sfd$ if and only if $f\in\Qfd$ and
$f^*:\Locfd(Y)\to\Locfd(X)$ is essentially surjective.
\item Equivalently, $f\in\Sfd$ if and only if $\pi_0f$ is bijective,
the induced map
\[
 \pi_1(X,x)^{\alg,k}\longrightarrow\pi_1(Y,f(x))^{\alg,k}
 \tag{D.2}
\]
is an isomorphism of affine group schemes for every source point $x$,
and (D.0) holds. Here the algebraic completions and their displayed map
are defined over $k$.
\end{enumerate}
The same equivalence test can be made on the ordinary homotopy categories
of the stable categories in (D.1). The ordinary category
$D^b(\Locfd(X))$ is not a replacement for $\mathcal P_k(X)$.
\end{claim}
\begin{proof}
First, (D.0) applied to the constant coefficient $k$ forces a bijection
on components. Indeed, the pullback $k^{\pi_0Y}\to k^{\pi_0X}$ is
injective only if every target component is hit, and is surjective only
if no two source components are identified; indicator functions prove
both assertions. Coefficients supported on one target component reduce
the remaining assertions to connected spaces.

For $M,N\in\Locfd(Y)$, the mapping spectrum in $\mathcal P_k(Y)$ is
the Eilenberg--Mac Lane spectrum of
$R\Gamma(Y;\underline{\operatorname{Hom}}_k(M,N))$.
Its comparison under $f^*$ is exactly the cohomology comparison for the
finite-dimensional local system $M^\vee\otimes_kN$. Thus (D.0) makes
all these mapping spectra equivalent. Closure under shifts, finite
cofibres, and retracts extends this assertion separately in both
variables. Each perfect local system has a finite Postnikov tower on a
connected component: its homotopy groups are finite-dimensional local
systems, and its fibre has finite amplitude. This proves full
faithfulness on the entire category. Conversely, full faithfulness
tested on $k$ and $L[n]$ gives (D.0). For infinitely many components,
mapping spaces are products of the componentwise mapping spaces; the
already proved component bijection proves the same assertion.

If all finite-dimensional local systems on $X$ come from $Y$, full
faithfulness lifts each attaching morphism of a finite Postnikov tower
and hence lifts its successive cofibres. It follows that every perfect
local system on $X$ comes from $Y$. Conversely, suppose $f^*$ is an
equivalence and $f^*E\simeq L$ with $L$ in degree zero. The component
bijection implies that $f$ meets every target component. Pointwise
cohomology then shows that $E$ has no cohomology outside degree zero,
so $E$ itself is an ordinary finite-dimensional local system.

The resulting equivalence of local-system categories is an equivalence
of $k$-linear tensor categories compatible with evaluation at the
specified basepoints. Such a category with its fibre functor recovers
the affine completion: its coordinate coalgebra is the coend of
$V^*\otimes V$ over finite-dimensional representations, modulo the
relations induced by intertwiners, and tensor product supplies the Hopf
multiplication. Restriction along (D.2) induces this exact coend map.
An equivalence preserving the fibre functor therefore makes (D.2) an
isomorphism. Conversely an isomorphism of affine completions identifies
their finite-dimensional representation categories with the specified
restriction. This proves the third description. Basepoint change
conjugates (D.2), which does not alter whether it is an isomorphism.

An exact functor between stable $\infty$-categories is fully faithful
when its maps on all shifted morphism groups are isomorphisms; those
groups are the homotopy groups of its mapping spectra. Essential
surjectivity does not change on passage to the homotopy category. This
justifies the ordinary homotopy-category formulation.

Finally, for $X=S^2$, $\Locfd(X)\simeq\operatorname{Vect}_{\mathrm{fd}}(k)$,
so $D^b(\Locfd(X))\simeq D^b(k)$ just as for a point. In contrast,
\[
 \operatorname{Hom}_{\operatorname{Ho}\mathcal P_k(S^2)}(k,k[2])
       =H^2(S^2;k)=k,
\]
whereas the corresponding morphism group at a point is zero. This also
rechecks R11: degree-zero local systems alone miss higher cohomology.
\end{proof}

\begin{claim}{D.3}{Proved}{The repaired class is a fundamental localizer on all small categories}
\label{D3}
For every field $k$, the class $\Sfd$, applied to nerve realizations,
satisfies identities, two-out-of-three, XVI Loc (1)c), retract closure,
Loc (2), and the relative Theorem A, Loc (3 bis). This assertion includes
empty and disconnected categories and imposes no accessibility
hypothesis on a pre-existing localizer.
\end{claim}
\begin{proof}
Restrictions to $\mathcal P_k(-)$ are contravariantly functorial.
Equivalences satisfy identities and two-out-of-three. If $gf$ and $fg$
are equivalences, $f$ has both a left and a right inverse, obtained by
composing $g$ with the inverses of these two composites; they agree.
This verifies Loc (1)c). Retracts of equivalences are equivalences.
A category with a final object has contractible realization, which
gives Loc (2).

We spell out relative Theorem A. Let $u:A\to B$, $q:B\to C$, and
$p=qu$ be functors. For $c\in C$ let $A/c=(p\downarrow c)$ and
$B/c=(q\downarrow c)$, and suppose $u/c:A/c\to B/c$ belongs to
$\Sfd$ for every $c$. Postcomposition in $C$ defines covariant
$\Cat$-valued functors $c\mapsto A/c$ and $c\mapsto B/c$.
Their Grothendieck constructions are denoted $E_A$ and $E_B$.
An object of $E_A$ is $(c,a,\alpha:p(a)\to c)$; an arrow to
$(c',a',\alpha')$ is a pair $(\beta:c\to c',\gamma:a\to a')$
such that $\alpha' p(\gamma)=\beta\alpha$.
The projection $r_A:E_A\to A$ sends this object to $a$.
The functor $j_A:A\to E_A$ sends $a$ to $(p(a),a,\mathrm{id})$.
The displayed arrow formula gives $j_A\dashv r_A$ and $r_Aj_A=\mathrm{id}$.
Consequently $Br_A$ and $Bj_A$ are mutually inverse weak homotopy
equivalences; the same holds for $B$.

The natural map $E_u:E_A\to E_B$ sends
$(c,a,\alpha)$ to $(c,u(a),\alpha)$ and satisfies
$r_BE_u=ur_A$ in $\Cat$. The Grothendieck-construction comparison gives
natural weak equivalences
\[
 \operatorname*{hocolim}_{c\in C}B(A/c)\longrightarrow BE_A,
 \qquad
 \operatorname*{hocolim}_{c\in C}B(B/c)\longrightarrow BE_B.
 \tag{D.3}
\]
The version and variance used here are the covariant form of
\citet[Theorem 26.8 and Corollary 26.9]{AChacholskiScherer},
whose Grothendieck construction in section 38 has arrows
$(\beta,\gamma)$ with the same covariance as our displayed formula.
In the simplicial
replacement, its map sends a chain $c_0\to\cdots\to c_n$ and a
chain in the initial fibre to the chain transported through the
successive fibres in the total category. Thus the comparison is natural
in $u$, not an abstract isomorphism of homotopy types.

For any small diagram of spaces $Z_c$ and any $\infty$-category
$\mathcal E$, restriction gives
\[
 \Fun(\operatorname*{hocolim}_c Z_c,\mathcal E)
       \simeq\operatorname*{holim}_{c\in C^{\mathrm{op}}}
                  \Fun(Z_c,\mathcal E).
 \tag{D.4}
\]
One can verify (D.4) at every level of the complete Segal nerve: its
$m$th space is $\operatorname{Map}(Z_c,
\Fun([m],\mathcal E)^{\simeq})$, and mapping out of a homotopy colimit
is a homotopy limit. Applying (D.4) with $\mathcal E=\Perf(k)$ makes
the map induced by (D.3) an equivalence, since all the fibre
restrictions are equivalences. Hence $E_u\in\Sfd$; the strictly
commuting projection square and two-out-of-three give $u\in\Sfd$.
This proves the relative axiom without an assumption that $\Perf(k)$
has arbitrary limits or defines a derivator. If $C$ is empty, $A$ and
$B$ are both empty. If just a comma category is empty, its functor
category is terminal and the same calculation applies.
\end{proof}

\subsection{What To\"en's theorem identifies}

\begin{claim}{D.4}{Known, with derivation}{The precise schematization equivalence criterion}
\label{D4}
For every field $k$ and every pointed map $f:(X,x)\to(Y,y)$ of
connected small spaces,
\[
 f\in\Sfd\quad\Longleftrightarrow\quad
 (f\otimes k)^{\sch}:(X\otimes k)^{\sch}\longrightarrow
                         (Y\otimes k)^{\sch}
 \text{ is an equivalence of pointed stacks}.
 \tag{D.5}
\]
The class on small categories is independent of chosen points and is
obtained by requiring a component bijection and (D.5) on every
component. The theorem requires neither algebraic closure, nor
characteristic zero, nor finite type or nilpotence of the spaces.
\end{claim}
\begin{proof}
We use the attached source of \citet{D_Toen2006}, checked against the
author's PDF. Definition 3.1.1 defines $P$-equivalences of connected
pointed simplicial presheaves by mapping-space equivalences into
$K(G,V,n)$ for affine finite-type $G$, finite-dimensional algebraic
$G$-representations $V$, and $n>0$. Lemma 3.3.2 identifies the induced
test on spaces with exactly the two conditions of D.2(3) on a
connected component. Its proof identifies the mapping spaces with
those into $K(G(k),V(k),n)$; Corollary 3.3.3 records that comparison.
Thus it applies over every field, including nonclosed fields and
positive characteristic.

Theorem 3.3.4 constructs a $P$-equivalence
$\eta_X:X\to(X\otimes k)^{\sch}$ with $P$-local target; its construction
uses a degreewise free simplicial-group resolution of $X$, affine
completion, and affination. Since both $\eta_X$ and $\eta_Y$ are
$P$-equivalences, (D.5) follows from two-out-of-three and the fact that
a $P$-equivalence between $P$-local objects is an equivalence. For the
latter fact, testing against each of the two local objects gives
mutually inverse morphisms in the homotopy category, so this is a
specified-map assertion. Corollary 3.3.6 identifies the resulting
functor as left adjoint to the pointed connected component of derived
global sections; Corollary 3.3.7 identifies its first homotopy sheaf
with $\pi_1(X,x)^{\alg,k}$.

For comparison, the attached
\citet{D_KatzarkovPantevToen2008} gives Definition 1.1.6,
Theorem 1.1.7, and Proposition 1.1.8 for the definition, existence,
and first homotopy sheaf, respectively, over $\mathbb C$.
Its introduction states separately both the representation-category
and finite-dimensional local-cohomology properties. It does not
discard the first condition of To\"en's Lemma 3.3.2. This is the exact
point missed by Lead 2.

For an unpointed connected map we may choose a source point and use
its image as target point. Different choices conjugate the
fundamental-group maps and canonically transport the local-system
tests, leaving (D.5) invariant. For arbitrary components we use the
groupoid formulation in D.2. Requiring only a test at one point of a
disconnected space would be insufficient: the inclusion of one point
into a two-point discrete space passes that component's test and fails
$H^0(-;k)$. The empty category likewise cannot be silently assigned
a basepoint or identified with a point.
\end{proof}

The word ``localization'' in (D.5) has a precise safe meaning:
schematization inverts exactly these maps, and its unit is a
$P$-localization \emph{in simplicial presheaves on affine schemes}.
Corollary 3.3.6 alone does not identify the localization of spaces with
the category of all schematic homotopy types. D.8 below gives a
concrete failure of a stronger interpretation using global sections.

\subsection{Uniform strictness examples}

\begin{claim}{D.5}{Proved}{An explicit finite presentation invisible to every field's finite representations}
\label{D5}
Let $K$ be the presentation CW complex with one vertex, four oriented
one-cells $a,b,c,d$, and four two-cells with attaching words
\[
 a^{-1}ba b^{-2},\quad b^{-1}cb c^{-2},\quad
 c^{-1}dc d^{-2},\quad d^{-1}ad a^{-2}.
 \tag{D.6}
\]
Then $K$ is connected and integrally acyclic, $\pi_1K\ne1$, and
every finite-dimensional representation of $\pi_1K$ over every field
is trivial. Consequently both $K\to *$ and its vertex inclusion
$*\to K$ belong to $\Sfd$ for every field.
\end{claim}
\begin{proof}
The cellular boundary $C_2(K;\mathbb Z)\to C_1(K;\mathbb Z)$ is
\[
 \begin{pmatrix}
 0&0&0&-1\\ -1&0&0&0\\0&-1&0&0\\0&0&-1&0
 \end{pmatrix},
 \qquad\det=-1.
 \tag{D.7}
\]
It follows directly that $\widetilde H_*(K;\mathbb Z)=0$.
We do not need to assert that $K$ is aspherical.

For completeness we check the group-theoretic input. Write
$G_{ab}=\langle a,b\mid a^{-1}ba=b^2\rangle$ and define
$G_{bc},G_{cd},G_{da}$ cyclically. In each such Baumslag--Solitar
group both indicated cyclic subgroups are infinite. For example,
$a(t)=t/2$, $b(t)=t+1$ define affine transformations of $\mathbb Q$
satisfying the relation, so $b$ is infinite; sending $a$ to $1$ and
$b$ to $0$ in $\mathbb Z$ proves that $a$ is infinite.
In $G_{ab}*_{\langle b\rangle}G_{bc}$, normal forms show that
$\langle a,c\rangle$ is a free group on $a,c$: nonzero powers of $a$
do not lie in the amalgamated cyclic subgroup by abelianization of
$G_{ab}$, and nonzero powers of $c$ do not lie there by abelianization
of $G_{bc}$ and infinitude of $c$. The same argument gives a free
$\langle c,a\rangle$ in $G_{cd}*_{\langle d\rangle}G_{da}$.
Amalgamating these two groups along that free group has precisely
presentation (D.6). The normal-form theorem embeds both factors,
so the resulting group $H=\pi_1K$ contains the infinite element $a$.
This is the decomposition recorded in the introduction of
\citet{D_Baker2019}; we have verified the embedding conditions here.

To show that $H$ has no nontrivial finite quotient, let the four
generator images have finite orders $m_i$. Each relation makes a
generator conjugate to its square, so all $m_i$ are odd. If some
$m_i>1$, let $p$ be the smallest prime dividing any $m_i$, and choose
$i$ with $p\mid m_i$. Conjugation by the preceding generator acts
as squaring on the cyclic subgroup generated by the $i$th generator.
Its order modulo $p$ is $\operatorname{ord}_p(2)>1$ and divides
$m_{i-1}$. But every prime divisor of $\operatorname{ord}_p(2)$ is
strictly smaller than $p$, since that order divides $p-1$. This
contradicts the choice of $p$. Thus every generator image is trivial.

A finitely generated subgroup of $\operatorname{GL}_n(k)$ is
residually finite over \emph{any} field. Here is the exact application
of the opened proof in \citet[Lemma 3.2(iii),(iv), and the proof of
Theorem 3.1]{D_Nica2013}. All entries of generators and their inverses
lie in a finitely generated integral domain $R\subset k$. For a
nonidentity matrix choose a nonzero entry of its difference from the
identity. The intersection of maximal ideals of $R$ is zero, and
every maximal residue field is finite. Reducing modulo a maximal
ideal avoiding that entry gives a finite matrix quotient in which
the matrix remains nonidentity. Therefore a nontrivial representation
of $H$ would give a nontrivial finite quotient of $H$, which is
impossible.

Every finite-dimensional local system on $K$ is now constant.
The finite free cellular chain complex (D.7) remains acyclic after
tensoring with $k$, and dualizing over a field gives its cohomology
version. Thus (D.0) holds for $K\to *$. Its representation category
is exactly that of a point, so D.2 makes this map an $\Sfd$-equivalence;
its section belongs to $\Sfd$ by two-out-of-three. The program checks
the relator exponent sums, determinant, and affine transformation
identity. The no-finite-quotient and residual-finiteness arguments
above supply the unbounded proof; no finite search is substituted.

To obtain literal small-category examples, let
$A_K=\Delta/\operatorname{Sing}K$. Its objects are singular simplices
$\sigma:\Delta^n\to K$; a morphism $\tau\to\sigma$ is a monotone map
$\alpha$ with $\tau=\sigma\alpha$. The canonical realization comparison
$BA_K\to K$ is a weak equivalence, natural in maps of spaces, and the
vertex supplies a functor $e\to A_K$. Restriction along this comparison
is an equivalence of every coefficient category used here. Thus all
membership and nonmembership statements transfer in both directions.
\end{proof}

\begin{claim}{D.6}{Proved}{Comparison with the two GALOIS classes}
\label{D6}
For every field $k$, the following are strict inclusions of classes of
functors of small categories:
\[
 W^{\mathrm{lc}}_{D(k)}\subsetneq\Sfd\subsetneq\Qfd\subsetneq W_{D(k)},
 \qquad
 W^{\mathrm{lc}}_{D(k)}\subsetneq W^1_{D(k)}\subsetneq\Qfd.
 \tag{D.8}
\]
The classes $\Sfd$ and $W^1_{D(k)}$ are incomparable. Over $\mathbb Q$,
$W^{\mathrm{lc}}_{D(\mathbb Q)}$ retains the actual discrete fundamental
group and is the class of $\pi_1$-rational equivalences; $\Sfd$ retains
only its affine completion together with the local cohomology test.
\end{claim}
\begin{proof}
We recheck the parts of GALOIS C.2, C.7, C.13, and C.14 that we use.
Writing $\mathcal M_k(X)=\Fun(X,D(k))$, its restriction is an
equivalence precisely when $f\in W^{\mathrm{lc}}_{D(k)}$, and is
fully faithful precisely when $f\in W^1_{D(k)}$. Evaluation at every
source point commutes with restriction. An equivalence forces a
component bijection; a perfect target fibre is then characterized by
its pullback being perfect, so an equivalence on $\mathcal M_k$
restricts to an equivalence on $\mathcal P_k$. This proves
$W^{\mathrm{lc}}_{D(k)}\subseteq\Sfd$.
Full faithfulness on all derived local systems gives full faithfulness
on its perfect subcategory, so $W^1_{D(k)}\subseteq\Qfd$.
Testing (D.0) on $k$ gives ordinary $k$-cohomology equivalence, which
over a field is the same as $k$-homology equivalence: the dual of a
linear map is invertible if and only if the map is invertible.
These observations prove all non-strict inclusions in (D.8).

For the necessary nonmembership tests, the left Kan extension counit
at $y$ is
\[
 C_*(F_y;k)\otimes_k N_y\longrightarrow N_y,
 \qquad F_y=X\times_Y^h\{y\}.
 \tag{D.9}
\]
The comparison comes from paths to $y$, which identify the restricted
coefficient diagram with the constant fibre $N_y$. Thus full
faithfulness is equivalent to $F_y$ being augmented $k$-acyclic for
every $y$, confirming exactly the $W^1$ test in GALOIS C.7.
On a connected component, evaluation is monadic with monad
$C_*(\Omega X;k)\otimes_k(-)$: its left adjoint is induction from
the point, and evaluation is conservative and preserves colimits.
An equivalence of restriction forces its extension--restriction unit
on the free module to be the quasi-isomorphism
$C_*(\Omega X;k)\to C_*(\Omega Y;k)$.
On $H_0$ this is the specified map
$k[\pi_1X]\to k[\pi_1Y]$, and bijectivity forces the actual group map
to be an isomorphism. This is the precise part of C.7 needed below.

Now use D.1 and D.5. The map $BC_\ell\to *$ belongs to $W^1_{D(k)}$
by (D.9), but not to $\Sfd$, because the regular representation is
not trivial. Its section $*\to BC_\ell$ belongs to $W_{D(k)}$ but
not to $\Qfd$. The map $K\to *$ belongs to $\Sfd$ and $W^1_{D(k)}$
but not to $W^{\mathrm{lc}}_{D(k)}$, since its fundamental group is
nontrivial. Its section $*\to K$ belongs to $\Sfd$ and therefore
$\Qfd$, but not to $W^1_{D(k)}$: its homotopy fibre is $\Omega K$,
whose component set is the nontrivial group $H$, so its augmented
$H_0(-;k)$ is not $k$. These four maps prove every asserted
strictness and both noninclusions.

Finally, for $k=\mathbb Q$, the loop-chain criterion above is the
cobar criterion in \citet[Theorem 3]{D_RiveraWierstraZeinalian2021},
which is equivalent to isomorphism on $\pi_1$ and on
$\pi_n\otimes\mathbb Q$ for $n\geq2$. This confirms GALOIS C.13.
The inclusion in GALOIS C.14 follows from the proved perfect-category
restriction, and D.5 gives a uniform finite-presentation strictness
example in place of its cardinal-controlled group. We do not identify
a rational homology equivalence with a $\pi_1$-rational equivalence.
\end{proof}

\begin{claim}{D.7}{Refuted}{Schematization equivalences satisfy none of the reverse excision axioms}
\label{D7}
For every field $k$, the fundamental localizer $\Sfd$ fails XVI
Loc (6$'$b), XII W(7 bis), cote 105 W7 bis, and both cube versions
XII W(7) and cote 105 W8 bis. The failure already occurs in a
three-object category with six arrows when $\operatorname{char}k\ne2$
and seven arrows when $\operatorname{char}k=2$.
\end{claim}
\begin{proof}
Use $\ell$ as in D.1. Define $X=\mathbf B C_\ell\star\{t,s\}$:
there are objects $b,t,s$, endomorphism group $C_\ell$ at $b$,
only the identity at $t$ and $s$, exactly one arrow $b\to t$ and
one $b\to s$, and no other arrows. Precomposition satisfies $u\circ g=u$ for each $u:b\to t$ or $u:b\to s$ and every $g\in C_\ell$.
The full subcategories $U=\{b,t\}$ and $V=\{b,s\}$ are sieves,
cover $X$, and intersect in $\mathbf B C_\ell$.
Their nerves are the two simplicial cones on $N\mathbf B C_\ell$,
so $BX$ is the unreduced suspension $\Sigma BC_\ell$.
The two inclusions of the base nerve are monomorphisms, so their
ordinary union is also their homotopy pushout; this identifies the
actual open-sieve square with the suspension square.

The space $BC_\ell$ is connected, and hence its suspension is simply
connected. The cone Mayer--Vietoris sequence and D.1 give
$\widetilde H_*(\Sigma BC_\ell;k)=0$. Every finite-dimensional local
system on it is constant, so $BX\to *$ belongs to $\Sfd$ by D.2.
The category $U$ has terminal object $t$, so $U\to X$ belongs to
$\Sfd$ by two-out-of-three. But $U\cap V\to V$ does not: it induces
$C_\ell^{\alg,k}\to1$, and the regular representation proves the
source affine group nontrivial. This refutes the reverse square.

For the arbitrary-span form use $\mathbf B C_\ell\to e$ on both
legs. Its cofibred Grothendieck construction is exactly $X$ above.
For the fibred convention of cote 105 the total category instead has
one arrow $t\to b$ and one $s\to b$, with the same group at $b$.
Sending each group element to its inverse identifies it with $X^{\mathrm{op}}$.
Reversal of each geometric simplex gives a natural homeomorphism
$B(X^{\mathrm{op}})\cong BX$; on the group fibre it induces the
weak equivalence given by reversing paths. Thus it preserves both the
triviality test for the total and the nontrivial representation test
for the central fibre.
The two point fibres map by $\Sfd$-equivalences to the total
construction, while either base-to-point leg fails. This is a
failure of the actual two biconditionals, not just of a convenient
different square.

For W8 bis take the morphism from this span to
$(e\leftarrow e\to e)$. The two outer fibre maps are identities,
the total map $X\to(e\star\{t,s\})$ is in $\Sfd$ since both totals
are $\Sfd$-aspherical, and the central map
$\mathbf B C_\ell\to e$ is not. Thus ``total plus the two outer
maps imply the middle map'' fails. Equivalently, map the above
open-sieve square to the corresponding cover of
$e\star\{t,s\}$; this is a strict cube in $\Cat$ and refutes XII
W(7). The program checks all category axioms and sieve conditions.
\end{proof}

\subsection{Field hypotheses and the meaning of localization}

\begin{claim}{D.8}{Refuted}{Derived global sections need not realize a localization of spaces}
\label{D8}
The assertion that $X\to R\Gamma((X\otimes k)^{\sch})$ is always
an $\Sfd$-equivalence is false, even for an algebraically closed
characteristic-zero field and a simply connected finite-type space.
In particular the adjunction in To\"en's Corollary 3.3.6 is not,
without further hypotheses, a reflective localization of spaces.
\end{claim}
\begin{proof}
Take $k=\mathbb C$ and $X=K(\mathbb Z,2)$.
The calculation in the proof of To\"en's Theorem 2.5.1, together
with Corollary 3.3.8, gives
$(X\otimes k)^{\sch}\simeq K(\mathbb G_a,2)$.
The additive group is a quasi-coherent sheaf and is acyclic on an
affine base, so its derived global sections are $K(\mathbb C,2)$.
The unit induces the inclusion $\mathbb Z\to\mathbb C$ on $\pi_2$.
By the first nonzero-degree Hurewicz theorem, its map on degree-two
homology over $\mathbb C$ is
\[
 \mathbb C\longrightarrow\mathbb C\otimes_{\mathbb Q}\mathbb C,
 \qquad z\longmapsto1\otimes z.
 \tag{D.10}
\]
Choose $a\in\mathbb C\setminus\mathbb Q$ and a $\mathbb Q$-linear
functional $\lambda:\mathbb C\to\mathbb Q$ with
$\lambda(1)=0$, $\lambda(a)=1$. Applying
$\lambda\otimes\mathrm{id}$ kills the image of (D.10) but does not
kill $a\otimes1$. Hence (D.10) is not surjective, so the unit fails
even the constant-coefficient test. The same argument works for any
proper field extension $k/\mathbb Q$.

There is a narrower positive statement, explicitly different from
D.4: To\"en's Theorem 2.5.1 and Corollary 2.5.4 give a fully faithful
affination model of the homology-localized category on nilpotent
finite-type spaces for $k=\mathbb Q$, or for algebraically closed
$k$ of positive characteristic. For simply connected such spaces,
Corollary 3.3.8 identifies affination with schematization. None of
these additional hypotheses is required for the equivalence test
(D.5); they concern reconstructing a localized \emph{space} by
global sections. Thus they cannot repair the counterexample D.1.
\end{proof}

\begin{claim}{D.9}{Proved}{The hypotheses and omissions that matter}
\label{D9}
The corrected general statement is D.2--D.4: for every field and all
small categories, equivalence of coherent perfect local-system
categories is a fundamental localizer; on each pointed connected
component its maps are exactly those inverted by To\"en's
schematization. The proalgebraic fundamental-group condition and the
cohomology condition are separately necessary. Basepoints merely
express the condition and algebraic closure and characteristic zero
are unnecessary.
\end{claim}
\begin{proof}
Dropping the proalgebraic condition admits $BC_\ell\to *$ from D.1,
which changes the finite-dimensional representation category.
Dropping the higher-cohomology condition admits $S^2\to *$, which
has trivial fundamental groups and fails $H^2(-;k)$. Keeping a
single chosen component while dropping the component-bijection
condition admits $*\to\{0,1\}$ and misses its $H^0$ defect.
Replacing coherent derived local systems by the derived category of
ordinary local systems again admits $S^2\to *$, by D.2. These are
the failures of omitted data. In contrast there is no failure upon
dropping algebraic closure or characteristic zero, because D.4 was
proved without either assumption. For the zero ring, excluded from
the definition of a field, the perfect coefficient category is zero
and every map passes the coefficient test. We make no assertion that
To\"en's theorem over a field applies to the zero ring or to an
arbitrary nonfield base.
\end{proof}

\begin{claim}{D.10}{Proved}{What the manuscript passages do and do not assert}
\label{D10}
Cote 133/006--008 does not state the finite-dimensional-local-system
localizer conjecture. The statement in Part D.1 is an additional
question, and Lead 2's proposed localizer argument cannot prove its
unrepaired version.
\end{claim}
\begin{proof}
The supplied transcription at 133/006, lines 30--44, reduces
universal formula questions for exterior-power operations on Witt
groups to universal quadratic bundles on
$B_{O(m)\times O(n)}$ and $B_{O(m)}$. Scans 007--008 discuss the
hyperbolic map from a Grothendieck group of locally free modules,
the Witt quotient, and the possible failure of hyperbolicity in a
classifying topos when two is not invertible. These are not
assertions about equivalences of homotopy types detected by local
systems; we do not promote the handwritten marginal doubts into
proved Witt-group theorems.

In \emph{Pursuing Stacks}, section 110, pages 443--444, Grothendieck
discusses possible schematization by Postnikov stages, including a
version retaining the fundamental groupoid discretely, and records
the obstruction suggested by torsion and the universal-coefficient
formula. This is motivation, not the criterion of To\"en's Lemma
3.3.2. The present additions are the counterexample D.1, the exact
perfect-local-system repair D.2--D.4, and the comparisons and
excision failures D.6--D.7. A comma spectral sequence can support
a cohomology comparison; it cannot supply the missing
two-out-of-three implication refuted in D.1.

Finally, the sentence following Definition 3.1.1 in the supplied
\emph{Champs affines} source reverses an implication. The correct
direction is $P$-equivalence implies $\mathcal O$-equivalence,
as used again in Lemma 3.2.1 and Corollary 3.3.8 of that source.
The opposite direction is refuted by $BC_\ell\to *$ over a field
of characteristic different from $\ell$. This confirms the specific
textual correction already recorded in GALOIS C.14.
\end{proof}
